\chapter{Nurses}

\medskip
\noindent\textit{\textbf{Under review.} This chapter is an AI-assisted educational draft; it carries no source citations and is not a substitute for professional medical advice.}

\section*{Definition and Overview}
Nursing is a health profession focused on the care of individuals, families, and communities in the promotion of health, the prevention of illness, and the care of the sick, disabled, and dying. In Australia, registered nurses (RNs), enrolled nurses (ENs), and nurse practitioners (NPs) form the largest component of the health workforce. Nursing practice encompasses assessment, diagnosis, planning, implementation, and evaluation of care across hospital, community, aged-care, and primary-health settings. The profession is regulated nationally by the Nursing and Midwifery Board of Australia (NMBA) under the Health Practitioner Regulation National Law, with registration standards that define scope of practice. Modern nursing integrates clinical skill with evidence-based practice, patient advocacy, and increasingly advanced and autonomous roles.

\section*{Epidemiology}
Nursing is the largest regulated health profession in Australia. As of 2024, there are approximately 460,000 registered nurses and midwives holding current registration with the NMBA. Around 44\% of registered nurses work in hospitals, with growing numbers in primary care, aged care, and community settings. Approximately 10--15\% of the nursing workforce identifies as male; the profession remains female-dominated. Rural and remote Australia faces chronic nursing shortages, with Indigenous health services and aged care particularly understaffed. Nursing is one of the fastest-growing occupations nationally, with the Commonwealth projecting a shortfall of tens of thousands of nurses over the coming decade, driven by an ageing population and an ageing workforce (the median nurse age is in the mid-40s, with a large cohort approaching retirement).

\section*{Aetiology and Pathophysiology}
The term ``nursing workforce'' spans several distinct regulated roles defined by education and scope of practice, rather than a disease process:
\begin{itemize}
    \item \textbf{Enrolled Nurse (EN):} Diploma-level (1.5--2 years); works under the direction and supervision of an RN; medication administration and some care planning within delegated scope
    \item \textbf{Registered Nurse (RN):} Bachelor-level (3 years) or postgraduate; autonomous assessment, planning, delivery, and evaluation of care; medication administration without direction
    \item \textbf{Nurse Practitioner (NP):} Master's-level with extended prescribing, diagnostic, and referral authority; works in primary care, emergency, and specialty services
    \item \textbf{Registered midwife:} Provides antenatal, intrapartum, and postnatal care across the full scope of maternity
    \item \textbf{Specialty and advanced practice:} Clinical nurse specialist (CNS), clinical nurse consultant (CNC), nurse educator, nurse manager, nurse researcher
\end{itemize}
Workforce supply is shaped by training places, retention, burnout, attrition, and international migration rather than by disease mechanisms. Staffing shortfalls are driven by heavy workloads, shift work, workplace violence, and limited career progression.

\section*{Clinical Features}
The ``clinical picture'' of the nursing role is the scope of care delivered:
\begin{itemize}
    \item \textbf{Assessment:} comprehensive physical, psychosocial, functional, and risk assessment; vital signs; early warning scores (e.g., NEWS2, adult deterioration detection system / ADDS)
    \item \textbf{Planning and coordination:} developing individualised care plans, coordinating multidisciplinary care, discharge planning
    \item \textbf{Interventions:} medication administration, wound care, catheter and line care, intravenous therapy, patient education, life-support (CPR, defibrillation)
    \item \textbf{Surveillance:} early recognition of clinical deterioration and escalation to medical staff
    \item \textbf{Advocacy:} representing patient preferences, advance care directives, and safety concerns
    \item \textbf{Health promotion:} screening, vaccination, chronic-disease self-management education, lifestyle counselling
    \item \textbf{End-of-life care:} symptom management, palliation, support of families during dying and death
\end{itemize}

\section*{Differential Diagnosis}
Differential considerations relevant to nursing scope-of-practice and career pathways:
\begin{itemize}
    \item \textbf{Registered nurse vs enrolled nurse:} differ in education, autonomy, and whether medication can be administered without direction
    \item \textbf{Nurse vs midwife:} midwifery is a distinct registration with a defined maternity scope
    \item \textbf{Nurse vs allied health:} physiotherapy, occupational therapy, and pharmacy have separate boards and scopes
    \item \textbf{Nurse practitioner vs doctor:} NPs have prescribing authority but not the full medical scope of diagnosis/surgery
    \item \textbf{Assistant in Nursing (AIN)/personal care worker:} unregulated support role working under delegation, not NMBA-regulated
\end{itemize}

\section*{When to Seek Medical Care}
A person does not ``seek medical care'' for the nursing role itself. Rather, patients under nursing care should escalate any concerning symptoms to their nurse or primary provider, and nurses are bound to escalate clinical deterioration promptly.

\textbf{Call 000 for:}
\begin{itemize}
    \item Cardiac or respiratory arrest (unconscious, not breathing, no pulse)
    \item Severe breathing difficulty or cyanosis
    \item Chest pain suggestive of acute coronary syndrome or stroke symptoms
    \item Severe allergic reaction (anaphylaxis), uncontrolled bleeding, or seizures
    \item Any sudden deterioration in level of consciousness
\end{itemize}

\section*{Diagnosis and Investigations}
Assessment of nursing care quality and workforce planning uses the following:
\begin{itemize}
    \item \textbf{Registration checks:} NMBA public register verifies qualifications and any conditions on registration
    \item \textbf{Clinical assessment:} nursing process (assess, diagnose, plan, implement, evaluate); risk-screening tools (falls, pressure injury, delirium, malnutrition, venous thromboembolism)
    \item \textbf{Observation tools:} early warning scores, fluid-balance charts, Glasgow Coma Scale, pain scores
    \item \textbf{Workforce metrics:} nurse-to-patient ratios, staffing mix, retention and attrition rates, vacancy rates, overtime
    \item \textbf{Outcome indicators:} pressure injury rates, medication-error rates, hospital-acquired infection rates, patient satisfaction, unplanned readmission
\end{itemize}

\section*{Management}
Management of the nursing workforce and delivery of nursing care:
\begin{itemize}
    \item \textbf{Education and credentialing:} NMBA-approved programs, mandatory professional development, annual registration and recency-of-practice requirements
    \item \textbf{Staffing:} legislated nurse-to-patient ratios in some states (e.g., Victoria public hospitals); safe staffing frameworks nationally
    \item \textbf{Retention strategies:} flexible rostering, pay and conditions improvements, career pathways, clinical supervision, wellness and burnout programs
    \item \textbf{Patient safety:} open disclosure, incident reporting (e.g., Australian Institute of Health and Welfare / state health incident systems), clinical handover frameworks (ISBAR)
    \item \textbf{Standard precautions:} hand hygiene (WHO Five Moments), PPE, aseptic technique, immunisation of staff
    \item \textbf{Delegation and supervision:} clear delegation of EN/AIN tasks consistent with scope-of-practice standards
    \item \textbf{Advanced practice:} NP-led clinics in primary care and rural areas expanding access to care
\end{itemize}

\section*{Complications}
Potential adverse outcomes relevant to nursing practice:
\begin{itemize}
    \item \textbf{Clinical:} medication errors, pressure injuries, falls, hospital-acquired infections, missed deterioration leading to cardiac arrest
    \item \textbf{Workforce:} burnout, compassion fatigue, workplace injury (manual handling, sharps), occupational violence and aggression
    \item \textbf{Systemic:} understaffing contributing to nurse-to-patient ratio breaches, unsafe workloads, and reduced quality of care
    \item \textbf{Personal:} shift-work-related sleep disturbance, mental health strain, high attrition from the profession
\end{itemize}

\section*{Prognosis and Outcomes}
Nursing is not a disease, so prognosis refers to workforce sustainability and care outcomes rather than patient survival. High-quality nursing care is consistently associated with lower mortality, fewer complications, shorter hospital stays, and higher patient satisfaction. Hospitals meeting safe nurse-to-patient ratios and skill mix show measurable reductions in adverse events, including fewer in-hospital deaths. However, the Australian nursing workforce faces an ageing demographic and projected shortages; without improved retention and training capacity, understaffing is expected to worsen over the next decade. Effective retention, workplace culture, and expansion of advanced-practice and primary-care roles are central to sustaining high-quality, accessible care for the Australian population.

\section*{Prevention and Self-Care}
\begin{itemize}
    \item Maintain current NMBA registration and recency-of-practice requirements
    \item Engage in continuous professional development and evidence-based practice
    \item Attend to personal wellbeing: sleep, nutrition, stress management, peer support, and professional counselling when needed
    \item Use safe manual-handling techniques and standard precautions to prevent injury and infection
    \item Follow vaccination schedules recommended for health workers (influenza, COVID-19, hepatitis B, pertussis)
    \item Report and manage occupational stress and workplace violence early
    \item Maintain awareness of open disclosure, incident reporting, and patient-safety culture
\end{itemize}

\section*{Sources and Further Reading}
The following reputable sources provide further information for healthcare students and patients seeking reliable, current advice about nursing in Australia.

\begin{itemize}
    \item Nursing and Midwifery Board of Australia (NMBA). \textit{Registration standards and code of conduct.} https://www.nursingmidwiferyboard.gov.au/
    \item Australian Health Practitioner Regulation Agency (Ahpra). \textit{Public register of nurses and midwives.} https://www.ahpra.gov.au/
    \item Australian Nursing and Midwifery Federation (ANMF). \textit{Professional practice and workplace resources.} https://www.anmf.org.au/
    \item Royal College of Nursing Australia (ACN). \textit{Nursing professional development and advocacy.} https://www.acn.edu.au/
    \item Australian Institute of Health and Welfare (AIHW). \textit{Health workforce data on nurses and midwives.} https://www.aihw.gov.au/
    \item World Health Organization. \textit{State of the World's Nursing.} https://www.who.int/
    \item Healthdirect Australia. \textit{Health workforce and nursing careers information.} https://www.healthdirect.gov.au/
\end{itemize}

\clearpage
